\section{Phased migration}
\label{sec:migration}

No phase jumps directly to unified. Each phase is additive, lands behind
the selector of Section~\ref{sec:coexistence} in \texttt{shadow} first,
and flips to \texttt{unified} only on the harness exit criterion. Phases
are ordered by blast radius: the lowest-risk new object first, the
highest-trust custody (staking) last.

\subsection*{P1 --- X locks as first-class objects}

Add \texttt{LockTx}, \texttt{UnlockTx}, and \texttt{SettleTx} to X as
new transaction types, plus the \texttt{XLock} output type
(Section~\ref{sec:xlock}) and the lock state in xvm's state. No module
uses them yet; this phase delivers the spine itself and its own tests
(lock $\to$ unlock-by-owner, lock $\to$ unlock-by-timeout, lock $\to$
settle-by-proof, double-consume rejection, conservation per asset). It
extends the existing \texttt{lux.VerifyTx} conservation path
(\texttt{vms/xvm/txs/executor/syntactic\_verifier.go:58}) to count
locked outputs. Risk: contained to X; no module behavior changes.

\subsection*{P2 --- DEX uses \texttt{lockID} for order funding}

D gains an order-funding path keyed by \texttt{lockID} and a
\texttt{SettleTx}-to-X settlement path, \emph{alongside} the
\texttt{available}/\texttt{locked} ledger
(\texttt{lx/dex/pkg/dchain/state.go:54}). Run in \texttt{shadow}: every
order funds and settles both ways; the harness diffs fills and balances
(Section~\ref{sec:coexistence}, axis 1) and both conservation invariants
(axis 2). This is the first module proven equal. The DEX matcher,
\texttt{settleFills}, and book are unchanged; only the funding/settlement
custody is dualized.

\subsection*{P3 --- bridge uses X locks}

B moves to \texttt{LockTx(purpose=bridge)} funding and
\texttt{RequireNFT} operator gating, settling returns via
\texttt{SettleTx}. The \texttt{nftfx} capability is \IMPL{}
(\texttt{vms/xvm/fxs/fx.go}); this phase wires it as the operator gate.
Bridges carry the most external-trust risk, so they migrate after the
internal DEX path is proven but before staking. Foreign-chain proof
verification routes through \texttt{VerifyChainCommitment}.

\subsection*{P4 --- P staking uses X stake-locks}

P moves stake to \texttt{LockTx(purpose=stake)} producing a
\texttt{stakeLockID}, with reward/unlock finalized by
\texttt{SettleTx}. P drops its parallel staking custody --- the
\texttt{StakeOuts} outputs
(\texttt{txs/add\_permissionless\_validator\_tx.go:50}) and the
return-and-mint on P's own state
(\texttt{txs/executor/proposal\_tx\_executor.go:283--325}). This phase
\textbf{must also resolve the asset-mint contradiction}: P's
\texttt{CreateAssetTx} (\texttt{txs/executor/create\_asset\_tx.go:61})
either moves to X or becomes a settled result; P retaining independent
mint authority is incompatible with the unified model. Staking is last
because it is the highest-value, longest-lived custody and the one whose
divergence would be most costly --- it flips to \texttt{unified} only
after the longest \texttt{shadow} soak.

\subsection*{P5 --- general plugin-settlement registry}

Generalize \texttt{targetChainID}/\texttt{targetModuleID} into a
registry so any computing chain can settle through X under explicit
rules. Each registry entry records:

\begin{itemize}[leftmargin=2em,nosep]
\item \texttt{chainID} / \texttt{moduleID} --- the settler identity that
  \texttt{LockTx} and \texttt{SettleTx} scope to.
\item \texttt{verificationMethod} --- which proof
  \texttt{VerifyChainCommitment} requires: P3Q attestation, starkfri
  STARK, or both, and the threshold.
\item \texttt{operatorNFTRequirement} --- the capability (if any) gating
  privileged actions, checked via \texttt{RequireNFT}.
\item \texttt{settlementRules} --- allowed purposes, allowed output
  forms (visible / shielded), expiry bounds.
\end{itemize}

\noindent With P5 the five primitives plus the registry are the complete
public surface: a new module integrates by registering, not by adding a
new custody mechanism. This is the ``one and only one way to settle''
end state.

\subsection*{Ordering rationale}

\begin{center}
\begin{tabular}{@{}llll@{}}
\toprule
Phase & Adds & Risk if it diverges & Soak \\
\midrule
P1 & lock primitives on X      & contained to X            & short \\
P2 & DEX lock-funding          & forked fills (internal)   & medium \\
P3 & bridge lock-funding       & external-chain trust      & medium \\
P4 & staking lock-funding      & highest-value custody     & long \\
P5 & plugin registry           & opens the surface         & per-module \\
\bottomrule
\end{tabular}
\end{center}
